\chapter{Pratica Operativa con PostgreSQL}
\label{app:postgres_pratica}

Questa appendice raccoglie i comandi essenziali per l'interazione diretta con il database tramite l'interfaccia a riga di comando \texttt{psql}. Risulta fondamentale per il debugging rapido e la configurazione manuale dell'ambiente di sviluppo backend su macchine locali (come Fedora Workstation).
\section{Accesso, Databases e Schemi}
\subsection{Gestione del Servizio e Accesso}

Per iniziare a lavorare con PostgreSQL, è necessario assicurarsi che il demone sia attivo in background sul sistema operativo.

\begin{lstlisting}[language=TerminalBash, title=Comandi Systemd (Fedora/Linux)]
# Avvio del servizio
sudo systemctl start postgresql

# Accesso all'interfaccia interattiva come utente postgres
sudo -u postgres psql
\end{lstlisting}

\begin{itemize}
    \item \texttt{sudo systemctl start postgresql}: significa "usa i privilegi di amministratore (\texttt{sudo}) per ordinare al gestore del sistema e dei servizi (\texttt{systemctl}, "system control") di avviare (\texttt{start}) il demone PostgreSQL, ovvero il server del database che gira in background".

    \item \texttt{sudo -u postgres psql}: significa "usa i privilegi di amministratore (\texttt{sudo}) per assumere l'identità dell'utente di sistema chiamato \texttt{postgres} (\texttt{-u postgres}) e, con quell'identità, avvia il programma client di utility terminale (\texttt{psql}) per stabilire una connessione interattiva con il server database".
\end{itemize}

\begin{tip}[Avvio Automatico]
Se utilizzi frequentemente il database durante lo studio, puoi abilitare l'avvio automatico del servizio all'accensione del sistema con il comando \texttt{sudo systemctl enable postgresql}.
\end{tip}

\begin{warningbox}[L'utente di sistema postgres]
Di default, PostgreSQL in ambiente Linux crea un utente di sistema chiamato \texttt{postgres}. L'autenticazione locale si basa spesso sul metodo \textit{peer}, il che significa che devi assumere temporaneamente l'identità di quell'utente di sistema (tramite \texttt{sudo -u}) per poter accedere come superuser al database senza fornire una password.
\end{warningbox}

\subsection{Visualizzare i Database Esistenti}

All'interno dell'istanza di PostgreSQL, è possibile elencare tutti i database creati seguendo due approcci differenti: l'utilizzo di un meta-comando specifico del client \texttt{psql} oppure l'esecuzione di una query SQL standard sul catalogo di sistema.

\paragraph{Metodo Rapido: Meta-comandi di \texttt{psql}}

Il modo più immediato per ottenere la lista dei database consiste nell'utilizzare il comando nativo del terminale interattivo (che non richiede il punto e virgola finale):

\begin{lstlisting}[language=bash, title=Meta-comando psql]
\l
# Oppure in forma estesa:
\list
\end{lstlisting}

Questo comando mostra un riepilogo tabellare contenente dettagli fondamentali quali il nome del database (\textit{Name}), l'utente proprietario (\textit{Owner}), la codifica dei caratteri (\textit{Encoding}) e le regole di localizzazione linguistica (\textit{Collate} e \textit{Ctype}).

\begin{dettagli}[Gestione del Paginatore]
Se l'elenco dei database è troppo lungo per la schermata del terminale, \texttt{psql} reindirizzerà l'output a un paginatore di sistema (solitamente l'utility \texttt{less}). Se l'interfaccia mostra il simbolo dei due punti (\texttt{:}) o la dicitura \texttt{END} in basso, dando l'impressione che il terminale sia bloccato, è sufficiente premere il tasto \texttt{q} sulla tastiera per interrompere la visualizzazione e tornare al prompt interattivo \texttt{postgres=\#}.
\end{dettagli}

\paragraph{Metodo Standard: SQL Puro}

Nel caso in cui si stia operando dall'interno di un'applicazione (ad esempio un backend Java tramite driver JDBC) o non si stia usando il client interattivo \texttt{psql}, è necessario interrogare direttamente le tabelle dizionario del database utilizzando l'SQL standard:

\begin{lstlisting}[language=PostgreSQL, title=Query sul Catalogo di Sistema]
SELECT datname FROM pg_database;
\end{lstlisting}

La tabella di sistema \texttt{pg\_database} memorizza i metadati globali relativi a ogni singolo database presente e gestito all'interno dell'intera istanza (\textit{cluster}) di PostgreSQL.

\subsection{Ispezione e Visualizzazione degli Schemi}

Dopo aver configurato nuovi schemi all'interno del database, è fondamentale poterne verificare l'esistenza e analizzarne i metadati. Questa operazione può essere eseguita rapidamente tramite le utility del client \texttt{psql} oppure in modo programmatico con l'SQL standard.

\paragraph{Metodo Rapido: Meta-comandi di \texttt{psql}}

Il comando più immediato per visualizzare gli schemi creati nel database corrente è:

\begin{lstlisting}[language=bash, title=Meta-comando per l'elenco degli schemi]
\dn
\end{lstlisting}

La sintassi \texttt{\textbackslash dn} sta per \textit{"display namespaces"}, un termine ereditato dall'architettura interna di PostgreSQL. L'output restituirà una tabella contenente il nome dello schema (\textit{Name}) e il rispettivo utente proprietario (\textit{Owner}).

Nel caso in cui sia necessario esaminare anche i diritti di accesso associati a ciascun contenitore (ad esempio per verificare quali utenti backend dispongano dei privilegi di lettura o scrittura) e le eventuali descrizioni, si utilizza la versione estesa del comando:

\begin{lstlisting}[language=bash, title=Ispezione estesa degli schemi]
\dn+
\end{lstlisting}

\paragraph{Metodo Standard: Interrogazione del Catalogo Globale}

Il meta-comando \texttt{\textbackslash dn} filtra l'output mostrando principalmente gli schemi definiti dall'utente. Tuttavia, ogni istanza di PostgreSQL racchiude al suo interno strutture invisibili all'ispezione ordinaria ma vitali per il motore SQL. 

Per ottenere la mappatura totale ed esaustiva di tutti gli schemi presenti, si interroga il dizionario dei dati:

\begin{lstlisting}[language=PostgreSQL, title=Query completa sulle strutture logiche]
SELECT schema_name, schema_owner FROM information_schema.schemata;
\end{lstlisting}

\begin{dettagli}[Gli schemi di sistema predefiniti]
Analizzando l'output della query sulla vista \texttt{information\_schema.schemata}, oltre allo schema \texttt{public} e a quelli personalizzati, emergeranno entità sistemiche cruciali:
\begin{itemize}
    \item \textbf{\texttt{pg\_catalog}:} il nucleo informativo del DBMS. Contiene le tabelle di sistema di PostgreSQL, le funzioni integrate, gli operatori e la definizione di tutti i tipi di dato interni.
    \item \textbf{\texttt{information\_schema}:} un'interfaccia standardizzata conforme alle specifiche ANSI SQL. È presente in modo analogo in molti altri RDBMS commerciali e open source, e serve a garantire la portabilità delle query di ispezione dei metadati tra database differenti.
\end{itemize}
\end{dettagli}

\subsection{Operare con Schemi Multipli}

Una volta definiti più schemi all'interno dello stesso database, sorge la necessità di istruire il motore SQL su quale "cartella logica" utilizzare per la creazione o l'interrogazione delle tabelle. Esistono due metodologie principali per operare in un ambiente multi-schema.

\paragraph{Strategia 1: Nomi Qualificati (Approccio Esplicito)}

Il metodo più rigoroso e indipendente dal contesto consiste nell'utilizzare il \textbf{nome qualificato} (\textit{Fully Qualified Name}). La sintassi prevede di anteporre il nome dello schema a quello della tabella, separandoli con un punto (\texttt{schema.tabella}). 

Questo approccio garantisce che la tabella venga creata o interrogata esattamente nel contenitore desiderato, a prescindere da come è configurata la sessione corrente.

\begin{lstlisting}[language=PostgreSQL, title=Creazione tabelle con Nomi Qualificati]
-- Creazione di tabelle in schemi differenti
CREATE TABLE schema1.utenti (
    id SERIAL PRIMARY KEY,
    username VARCHAR(50) NOT NULL
);

CREATE TABLE schema2.ordini (
    id SERIAL PRIMARY KEY,
    totale NUMERIC(10, 2) NOT NULL
);
\end{lstlisting}

\begin{warningbox}[Nota Tecnica su SERIAL vs IDENTITY]
In questi primi esempi esplorativi è stato utilizzato il tipo \texttt{SERIAL} per la generazione degli identificativi. È importante precisare che \texttt{SERIAL} non è un vero e proprio dominio, ma una comodità sintattica storica e proprietaria di PostgreSQL, oggi considerata \textit{legacy}. Per garantire la massima portabilità e aderenza allo standard ANSI SQL, nelle sezioni successive dedicate al DDL si abbandonerà \texttt{SERIAL} in favore della clausola standard \texttt{GENERATED ALWAYS AS IDENTITY}.
\end{warningbox}

\begin{bestpractice}[Infrastruttura come Codice (IaC)]
Nello sviluppo backend professionale, la struttura del database viene gestita tramite script automatizzati di "migrazione" (utilizzando librerie Java come Flyway o Liquibase). In questi script si preferisce sempre l'approccio esplicito (\texttt{schema.tabella}) per garantire che l'esecuzione sia deterministica e che le tabelle vengano generate nell'ambiente corretto, a prescindere dalle configurazioni del server di produzione.
\end{bestpractice}

\paragraph{Strategia 2: Modifica del \texttt{search\_path} (Approccio Implicito)}

Se è necessario eseguire una lunga sequenza di operazioni (es. script di popolamento massivo) all'interno di un unico schema, dichiarare ripetutamente il nome qualificato risulta ridondante. 

È possibile alterare la variabile \texttt{search\_path} per indicare a PostgreSQL di puntare automaticamente a uno schema specifico qualora non venga fornito un prefisso.

\begin{lstlisting}[language=PostgreSQL, title=Modifica del search\_path]
-- Sposta il focus esclusivo della sessione corrente su "schema3"
SET search_path TO schema3;

-- Questa tabella verra creata implicitamente dentro "schema3"
CREATE TABLE prodotti (
    id SERIAL PRIMARY KEY,
    nome_prodotto VARCHAR(100) NOT NULL
);
\end{lstlisting}

\begin{warningbox}[Ambito di visibilità]
L'istruzione \texttt{SET search\_path} ha effetto \textbf{esclusivamente} sulla sessione attiva in quel momento. Se la connessione al database cade o se si apre un nuovo terminale client, il percorso di ricerca tornerà ai valori predefiniti globali (generalmente lo schema utente seguito da \texttt{public}).
\end{warningbox}

\paragraph{Verifica delle Strutture Multi-Schema}

Per avere un quadro completo e verificare in quali schemi risiedano effettivamente le tabelle del database (indipendentemente dal \texttt{search\_path} attuale), è possibile utilizzare un meta-comando \texttt{psql} sfruttando l'asterisco come carattere jolly per ignorare il filtro sul percorso di ricerca:

\begin{lstlisting}[language=bash, title=Ispezione globale delle tabelle]
\dt *.*
\end{lstlisting}
L'output presenterà una colonna aggiuntiva indicante lo schema di appartenenza di ogni singola entità.

\paragraph{Verifica del \texttt{search\_path} Corrente}

Poiché il prompt interattivo del client \texttt{psql} mostra unicamente il database di connessione e non lo schema attivo, è responsabilità dello sviluppatore interrogare esplicitamente il motore per conoscere il percorso di ricerca attuale prima di eseguire operazioni DDL.

\begin{lstlisting}[language=PostgreSQL, title=Comandi di ispezione della sessione]
-- Mostra l'intera stringa del percorso di ricerca configurato
SHOW search_path;

-- Restituisce unicamente il primo schema valido del percorso
SELECT current_schema();
\end{lstlisting}

\begin{dettagli}[La differenza tra SHOW e SELECT]
Il comando \texttt{SHOW} è un'istruzione amministrativa di PostgreSQL utilizzata per ispezionare i parametri di configurazione in esecuzione a livello di sessione (come il fuso orario, l'encoding o, appunto, il \texttt{search\_path}). 
La funzione \texttt{current\_schema()}, invece, è valutata a runtime e restituisce un singolo valore testuale. Questa distinzione la rende fondamentale quando si scrivono procedure archiviate (Stored Procedures) o script applicativi in cui il codice deve interrogare dinamicamente il database per sapere in quale contenitore sta operando.
\end{dettagli}

\subsection{Riepilogo Operativo: Flusso di Lavoro in \texttt{psql}}

Questo paragrafo sintetizza il ciclo di lavoro standard all'interno del terminale interattivo. Partendo dall'accesso iniziale, ecco la sequenza dei comandi più diretti e comunemente impiegati nel lavoro quotidiano per navigare tra le strutture logiche in modo rapido.

\begin{enumerate}
    \item \textbf{Visualizzare i database:} l'ispezione immediata delle istanze disponibili sul server.
    \item \textbf{Entrare nel database:} il comando per stabilire la connessione col database di lavoro.
    \item \textbf{Visualizzare gli schemi:} la scansione dei contenitori logici (namespaces) definiti al suo interno.
    \item \textbf{Impostare il path:} l'istruzione per orientare il focus della sessione su uno specifico schema.
    \item \textbf{Cambiare path:} la sovrascrittura della variabile di sessione per passare a un altro modulo.
    \item \textbf{Verificare il path corrente:} il controllo diagnostico fondamentale prima di eseguire operazioni DDL.
    \item \textbf{Uscire dal database:} poiché non esiste un comando di spegnimento della singola connessione, la prassi consiste nello spostarsi sul database amministrativo predefinito (\texttt{postgres}), liberando così il database precedente.
\end{enumerate}

\begin{lstlisting}[language=PostgreSQL, title=Cheat Sheet: Navigazione di Base]
-- 1. Visualizzare i database
\l

-- 2. Entrare nel database di lavoro
\c nome_database

-- 3. Visualizzare gli schemi
\dn

-- 4. Impostare il path per lo schema in cui lavorare
SET search_path TO nome_schema;

-- 5. Cambiare il path per lavorare su uno schema differente
SET search_path TO nuovo_schema;

-- 6. Verificare il path corrente
SHOW search_path;

-- 7. Uscire dal database (spostandosi su quello di default)
\c postgres
\end{lstlisting}

\begin{warningbox}[Meta-comandi vs SQL Standard]
Come si evince dal riepilogo, nel lavoro quotidiano si alternano continuamente i meta-comandi del client (es. \texttt{\textbackslash c} o \texttt{\textbackslash dn}, che non richiedono il punto e virgola finale) e le istruzioni SQL standard (come \texttt{SET} e \texttt{SHOW}, che invece lo esigono tassativamente per essere eseguite dal motore).
\end{warningbox}

\section{Definizione e Manipolazione dei Dati (DDL e DML)}

Una volta configurato l'ambiente e posizionato il \texttt{search\_path} sullo schema desiderato, il flusso di lavoro si sposta sulla definizione e sull'utilizzo delle strutture fisiche.

\subsection{Esplorazione del Contenitore Logico}

Prima di definire nuove tabelle o manipolare dati esistenti, è ragionevole verificare lo stato attuale dello schema di lavoro. Il client \texttt{psql} offre meta-comandi specifici per ispezionare le entità senza dover interrogare manualmente i dizionari di sistema.

\begin{lstlisting}[language=TerminalBash, title=Meta-comandi di Ispezione delle Tabelle]
# Elenca tutte le tabelle visibili nel search_path corrente
\dt

# Elenca le tabelle di uno schema specifico (ignorando il search_path)
\dt nome_schema.*

# Mostra l'anatomia dettagliata di una specifica tabella
\d nome_tabella
\end{lstlisting}

\begin{dettagli}[L'output diagnostico di \texttt{\textbackslash d}]
Il comando \texttt{\textbackslash d} (senza la "t") è lo strumento diagnostico più prezioso durante la stesura del DDL. L'output generato disseziona la tabella mostrando:
\begin{itemize}
    \item L'elenco completo degli attributi (colonne).
    \item I domini associati (tipi di dato come \texttt{integer}, \texttt{character varying}).
    \item I modificatori strutturali (vincoli in-line come \texttt{NOT NULL} o regole di auto-incremento come \texttt{IDENTITY}).
    \item In calce alla struttura, l'elenco esplicito degli indici, dei vincoli di chiave primaria (\texttt{PRIMARY KEY}) ed esterna (\texttt{FOREIGN KEY}), inclusi i riferimenti incrociati ad altre tabelle.
\end{itemize}
\end{dettagli}

\subsection{Data Definition Language (DDL) Pratico}

La sintassi di creazione delle tabelle su PostgreSQL rispecchia fedelmente lo standard SQL. Per la generazione delle chiavi surrogate, in ambiente moderno si predilige il costrutto standard \texttt{IDENTITY} anziché il vecchio tipo proprietario \texttt{SERIAL}.

\begin{lstlisting}[language=PostgreSQL, title=DDL: Creazione di Tabelle Relazionate]
-- Creazione della tabella padre (Utenti)
CREATE TABLE utenti (
    id_utente INT GENERATED ALWAYS AS IDENTITY,
    username VARCHAR(50) NOT NULL,
    email VARCHAR(255) NOT NULL,
    data_registrazione TIMESTAMP DEFAULT CURRENT_TIMESTAMP,

    -- Vincoli out-of-line
    CONSTRAINT pk_utenti PRIMARY KEY (id_utente),
    CONSTRAINT uq_utenti_email UNIQUE (email),
    CONSTRAINT uq_utenti_username UNIQUE (username)
);

-- Creazione della tabella figlia (Ordini)
CREATE TABLE ordini (
    id_ordine INT GENERATED ALWAYS AS IDENTITY,
    utente_id INT NOT NULL,
    totale NUMERIC(10, 2) NOT NULL,
    stato VARCHAR(20) DEFAULT 'IN_ELABORAZIONE',

    -- Vincoli out-of-line
    CONSTRAINT pk_ordini PRIMARY KEY (id_ordine),
    CONSTRAINT chk_ordini_totale CHECK (totale > 0),
    CONSTRAINT fk_ordini_utenti FOREIGN KEY (utente_id) 
        REFERENCES utenti(id_utente)
        ON DELETE RESTRICT
        ON UPDATE CASCADE
);
\end{lstlisting}

\begin{tip}[Ispezione rapida dei vincoli]
Dopo aver eseguito il DDL, per verificare che PostgreSQL abbia recepito correttamente i vincoli (inclusi i nomi assegnati come \texttt{pk\_utenti} o \texttt{fk\_ordini\_utenti}), è sufficiente utilizzare il meta-comando \texttt{\textbackslash d nome\_tabella} (es. \texttt{\textbackslash d ordini}). Il client interattivo mostrerà l'elenco delle colonne, seguito da una sezione dettagliata denominata "Indexes", "Check constraints" e "Foreign-key constraints".
\end{tip}

\subsection{Data Manipulation Language (DML) di Base}

Una volta definita l'architettura relazionale, si procede con l'inserimento (\texttt{INSERT}) e l'interrogazione (\texttt{SELECT}) dei record per collaudare le logiche di business configurate sul motore di persistenza.

\begin{lstlisting}[language=PostgreSQL, title=DML: Inserimento e Interrogazione]
-- Popolamento della tabella padre 
-- (L'ID viene escluso dall'elenco perche' GENERATED ALWAYS)
INSERT INTO utenti (username, email)
VALUES ('rookie_dev', 'rookie@example.com'),
       ('java_master', 'master@example.com');

-- Popolamento della tabella figlia
-- (Assumiamo per ipotesi che i due utenti generati abbiano ID 1 e 2)
INSERT INTO ordini (utente_id, totale)
VALUES (1, 150.50),
       (1, 45.00),
       (2, 890.00);

-- Query di verifica (DQL) con JOIN esplorativa
SELECT u.username, o.totale, o.stato
FROM utenti u
JOIN ordini o ON u.id_utente = o.utente_id;
\end{lstlisting}

\begin{warningbox}[Test dei Vincoli (Destructive Testing)]
In fase di sviluppo e debugging, prima di collegare il backend Java, è fondamentale testare attivamente la robustezza del database inviando query volutamente errate. Si consiglia di tentare forzature come: inserire un \texttt{utente\_id} inesistente nella tabella figli, inserire un ordine con totale negativo (violazione del \texttt{CHECK}), o tentare di bypassare l'inserimento manuale dell'ID. L'obiettivo è verificare che il DBMS respinga correttamente le transazioni non valide.
\end{warningbox}